% =========================================================================
% p2_sec05_discussion.tex — Part II discussion & policy (verbatim sec06).
% =========================================================================
% =========================================================================
% sec06_discussion_policy.tex — Section VI of the sealed master
% sealed master paper_seni142_ieee_transactions_en.md; editorial injection pass 2026-10-17.
% =========================================================================
\section{Discussion and Policy Implications}
\label{sec:discussion}

\textbf{Deterministic resilience for hurricane-exposed islands.} Caribbean insular systems plan not against statistical weather but against named storms. A benchmark whose $N-1$ margins carry correlation-honest bootstrap envelopes, whose largest credible contingency is exercised explicitly, and whose state resolution is a solve time compatible with the 1--5~s EMS cycle (which also includes acquisition, state estimation, and bad-data processing) converts resilience from an aspiration into a scheduled computation. For a country that has twice in twelve months mobilized around nationwide interruptions \cite{ocseni}, \cite{diarioLibreHainamosa_2026}, and where post-collapse restoration historically exceeds 10 to 14 hours due to Ferranti overvoltages ($B_{\mathrm{shunt}}$) on unloaded lines tripping black-start units, the difference between hours-long legacy studies and 2.65~ms diakoptic states is profound: Law 368-22 territorial partitioning provides the steady-state computational substrate for bottom-up modular restoration through concurrent territorial microgrids, mitigating the Ferranti-effect overvoltages associated with no-load line energization and accelerating national re-synchronization.

\textbf{Planning law as a computational boundary.} The partition of this benchmark is not an optimization artifact: it is the national territorial statute \cite{ley368}. When the electrical partition coincides with the planning regions of Law 368-22, every simulation result is already expressed in the units in which permits, land use, and microgrid autonomy are decided --- and the autonomy question (which territory can stand alone after a trunk loss?) becomes a property of the model rather than a separate study. The renewable incentive framework \cite{ley57} and the national expansion planning cycle \cite{cne2022} acquire, in consequence, a common evidentiary substrate.

\textbf{Standards alignment and data sovereignty.} The benchmark speaks the standards it must interoperate with: CIM-consistent model interchange \cite{iec61970}, measurement-grade power-quality definitions \cite{iec61000}, \cite{ieee1459}, harmonic practice \cite{ieee519}, study-method recommended practice \cite{ieee399}, and inverter-based-resource interconnection \cite{ieee2800}, \cite{ieee1547}. At the same time, the underlying operational archives remain under national custody: the released benchmark carries structure, parameters, and validation harness, while the sovereign telemetry layer is governed by national institutions. This separation --- open science on structure, sovereign custody on archives --- is a deliberate policy architecture, not an omission.

\textbf{Capacity building.} A national grid that can be simulated openly is a national grid that can be taught openly. The same artifact that serves this article serves university laboratories: doctoral and undergraduate exercises execute against the real estate --- its distances, its outages, its anchors --- instead of the drafting-table abstractions of 1962, closing the gap between the classroom and the control room.

\textbf{Limitations.} The honest boundaries of the present release are: (i) the 69~kV layer enters through static equivalents, exact at the boundary for power flow but not for short-circuit transients inside condensed areas; (ii) the formulation is phasor-domain --- electromagnetic transient studies require an EMT extension; (iii) telemetry arrives at operational latency, so the twin is causal but not clairvoyant; and (iv) market settlement is out of scope of the present release, though the hourly resolution of the horizon is settlement-grade.

\textbf{Replication for other island systems.} Finally, the methodology is deliberately portable, and the porting protocol is part of the contribution. Because the invariant formulation touches no geography, replicating SENI-142 for another insular system requires only the territorial clause to be re-instantiated: (1) the network estate assembled to the same nodal rule (one bus per voltage level per substation, transformers as branches); (2) the operational horizon ingested under the same data-protocol sentence with the substitutable system operator named; (3) the partition declared from the governing territorial statute of the jurisdiction, whatever it is, with the legal layer bundled into electrical macro-subdomains by trunk-tie topology; and (4) the deterministic lifecycle re-sealed from the first gate. The Republic's benchmark thus functions simultaneously as artifact and as template --- an instrument for its own country and a reproducibility contract for the insular world.
